\documentclass[border=7bp]{standalone}
% Shared preamble for the report's editable vector figures.
\usepackage{fontspec}
\usepackage{xeCJK}
\setmainfont{Calibri}
\setCJKmainfont{SimSun}[AutoFakeBold=1.5]
\usepackage{amsmath}
\usepackage{tikz}
\usetikzlibrary{arrows.meta,positioning,calc,shapes.geometric,fit}
\definecolor{Ink}{HTML}{25313D}
\definecolor{Blue}{HTML}{315E80}
\definecolor{PaleBlue}{HTML}{E7EFF5}
\definecolor{Muted}{HTML}{626B74}
\definecolor{PaleGray}{HTML}{F0F1F3}
\tikzset{
  >=Stealth,
  every node/.style={font=\fontsize{10bp}{13bp}\selectfont,text=Ink},
  flow/.style={->,draw=Ink,line width=.65pt},
  frame/.style={rectangle,draw=Ink,line width=.65pt,align=center,
    minimum height=.85cm,text width=5cm,inner sep=5pt},
  emphasis/.style={frame,draw=Blue,fill=PaleBlue},
  supporting/.style={frame,draw=Muted,fill=PaleGray},
  annotation/.style={font=\fontsize{9bp}{12bp}\selectfont,text=Muted,align=center},
}

\begin{document}
\begin{tikzpicture}
\path (-7.3,0) rectangle (7.3,0);

\node[emphasis,text width=11.8cm] (new) at (0,0) {生成候选后继：计算 $g$ 与 $h$，先检查是否已在 Closed};
\node[supporting,text width=11.8cm] (closed) at (0,-1.3) {Closed：两种方案均用 HashSet，平均 $O(1)$；本实验未更换该结构};
\node[emphasis,text width=4.8cm] (list) at (-3.2,-2.7) {ListFrontier：有序列表};
\node[emphasis,text width=4.8cm] (pq) at (3.2,-2.7) {PqFrontier：优先队列 + 状态索引};
\node[frame,text width=4.8cm] (l1) at (-3.2,-4.1) {定位重复状态：逐项比较 $O(N)$};
\node[frame,text width=4.8cm] (p1) at (3.2,-4.1) {定位重复状态：HashMap 平均 $O(1)$};
\node[frame,text width=4.8cm] (l2) at (-3.2,-5.8) {新状态或更小 $g$：保持列表有序\\插入、删除可能搬移元素 $O(N)$\\较差重复路径直接丢弃};
\node[frame,text width=4.8cm] (p2) at (3.2,-5.8) {新状态：堆插入 $O(\log N)$\\更小 $g$：旧堆项线性删除 $O(N)$\\较差重复路径直接丢弃};
\node[frame,text width=4.8cm] (l3) at (-3.2,-7.6) {取最优：移除列表首项 $O(N)$};
\node[frame,text width=4.8cm] (p3) at (3.2,-7.6) {取最优：堆出队 $O(\log N)$\\同时删除状态索引};
\foreach \a/\b in {new/closed,list/l1,pq/p1,l1/l2,p1/p2,l2/l3,p2/p3}{\draw[flow](\a)--(\b);}
\draw[flow](closed.south)-|(list.north);\draw[flow](closed.south)-|(pq.north);
\node[annotation,text width=11cm] at (0,-8.7) {图示为操作代价。额外索引增加空间；替换仍为线性删除，不能据此虚写全部操作的固定加速比};

\end{tikzpicture}
\end{document}
